% !TeX root = ../../Book4.tex
% 纲要标题：## 第8章　投射消解与内射消解
% 纲要标签：b4:ch:07
% 纲要位置：计划纲要.md，第 3372 行。

\chapter{投射消解与内射消解}
\label{b4:ch:07}
\BookChapterNumber{b4:ch:07}

\begin{chapterabstract}
本章构造投射与内射消解，证明比较定理和马蹄引理，说明消解的提升与唯一性应在同伦意义下理解，并给出短消解、Koszul 消解和 bar 消解的模型。
\end{chapterabstract}

\begin{learninggoals}
\begin{itemize}
\item 区分增广复形与去掉增广项的消解复形，说明投射和内射的次数方向。
\item 从核或余核逐次构造消解，并指出一般交换范畴所需的存在性假设。
\item 构造比较映射及两个提升之间的同伦，解释消解独立性。
\item 用马蹄引理构造相容消解，辨别逐次分裂与复形分裂。
\item 识别短消解的适用条件，并核验 Koszul 与 bar 微分的正合性来源。
\end{itemize}
\end{learninggoals}

设整数 \(n\ge2\)。映射 \(\mathbb Z\xrightarrow{n}\mathbb Z\twoheadrightarrow\mathbb Z/n\mathbb Z\) 把商模的生成元和关系都留在一条正合列里。给定从中间的自由模 \(\mathbb Z\) 到另一个模 \(M\) 的同态，它何时能经过商模 \(\mathbb Z/n\mathbb Z\)？答案取决于前一项的像是否被送到零。从自由模出发可以先选择生成元的像，再由前一项检查关系。对一般模，关系组成的核还可能有自己的关系，因此须继续取自由满射，逐层记录。自由模之所以适合这种构造，不仅因为可以指定基元的像，还因为从它出发的态射能沿满射提升。后面的比较定理反复使用的是这一提升性质，不需要依赖某组基；因此将自由模推广为具备同样性质的投射对象，就得到投射消解。接下来还须证明：换一种选择得到的消解，仍会给出相同的同调计算。

\ProseParagraphBreak
本章使用\chapref{b4:ch:04}的交换范畴与蛇引理、\chapref{b4:ch:05}的同调长正合列，以及\chapref{b4:ch:06}的链同伦。模论中的投射、内射及其存在性可复习第二卷第7.1—7.2节；其中定理7.2.8已完整证明任意环的模范畴有足够多内射对象。抽象交换范畴的结论总要另外声明足够多投射或内射对象，不能只由交换范畴公理推得。

\ProseParagraphBreak
Weibel 的《An Introduction to Homological Algebra》\cite[Sections 2.2--2.3, Comparison Theorem 2.2.6, Horseshoe Lemma 2.2.8, Comparison Theorem 2.3.7]{Weibel1994HomologicalAlgebra}给出本章前三节的基本构造。该书使用链指标；本章为展示投射消解保留非负链指标 \(P_n\)，同时以 \(P^{-n}=P_n\) 明确它与本书上链约定的转换。

% !TeX root = ../../Book4.tex
% 纲要标题：### 8.1 消解的构造
% 纲要标签：b4:sec:0701
% 纲要位置：计划纲要.md，第 3374 行。

\section{消解的构造}
\label{b4:sec:0701}

生成元给出自由满射，关系组成其核；继续为核取生成元，便得到消解。自由模可以沿满射提升映射，后面的比较定理只需要这一性质，因而可以把自由模放宽为投射对象。相反方向的构造使用内射对象的延拓性质；这两个方向都依赖范畴中相应对象是否足够多。


% 纲要内容（仅作写作计划，尚非教材正文）：
%
% * 模的自由及投射消解
% * 内射消解
% * 足够多投射或内射对象的假设
%

\subsection{模的自由及投射消解}
\label{b4:subsec:0701:01}

一个模的生成元给出自由满射，生成元之间的关系则成为这个满射的核。若核还不是容易处理的模，就继续为核选择生成元。消解把这项重复操作写成一个正合复形；它所保留的不是某一组最优生成元，而是所有关系都被逐层记录这一事实。

\begin{definition}\label{b4:def:projective-resolution}
设 \(\mathcal A\) 是交换范畴，\(P\) 为其对象。若对每个满态射 \(e:X\to Y\) 及态射 \(f:P\to Y\)，都存在态射 \(\widetilde f:P\to X\)，使 \(e\widetilde f=f\)，则称 \(P\) 为\term[tousheduixiang]{投射对象}[projective object]。给定对象 \(M\in\mathcal A\) 及增广链复形
\[
\cdots\longrightarrow P_2\xrightarrow{\mathrm{d}_2}P_1
\xrightarrow{\mathrm{d}_1}P_0\xrightarrow{\varepsilon}M\longrightarrow0
\]
若此序列正合，且每个 \(P_n\)（\(n\ge0\)）都投射，则称它为 \(M\) 的\term[toushexiaojie]{投射消解}[projective resolution]。若 \(R\) 为含幺环、\(\mathcal A=R\text{-}\mathbf{Mod}\) 为幺性左 \(R\) 模范畴，且各 \(P_n\) 自由，则称为\term[ziyouxiaojie]{自由消解}[free resolution]。
\end{definition}

正合性说明这条序列确实记录了 \(M\) 及各层关系；各项的投射性则保证从它出发的映射可以沿满态射逐层提升。后一条件将在\secref{b4:sec:0702}的比较定理中使用，不是正合性的另一个说法。

\ProseParagraphBreak
增广项 \(M\) 不属于非负次数的复形 \(P_\bullet\)：它的零次项是 \(P_0\)，从 \(P_0\) 向更低次数的微分为零。由增广诱导的商同构是
\[
H_0(P_\bullet)=P_0/\operatorname{im}\mathrm{d}_1
=P_0/\ker\varepsilon\xrightarrow{\;\overline\varepsilon\;}M.
\]
它是同构，因为 \(\varepsilon\) 满，且核正好是 \(\operatorname{im}\mathrm{d}_1\)；正次数同调则由正合性全部为零。写成上链复形时约定 \(P^{-n}=P_n\)，并在正次数补零，得到
\[
P^\bullet:\quad\cdots\longrightarrow P^{-2}\longrightarrow P^{-1}
\longrightarrow P^0\longrightarrow0.
\]
与它比较的是 \(M[0]\)，其中 \(M[0]^0=M\)，其余次数为零。增广给出复形映射 \(\varepsilon:P^\bullet\to M[0]\)：零次分量为原来的 \(\varepsilon\)，其余分量为零；\(\varepsilon\mathrm{d}^{-1}=0\) 保证复形映射条件。由刚才的同调计算，它是拟同构。

\begin{proposition}\label{b4:prop:free-resolution-exists}
设 \(R\) 是含幺环。每个幺性左 \(R\) 模都有自由消解，因而也有投射消解。每个幺性右 \(R\) 模也有相应的右模消解。
\end{proposition}
\begin{proof}
设 \(M\) 是幺性左 \(R\) 模，令 \(K_0=M\)。以 \(K_0\) 的底集合为基取自由模 \(P_0\)，把基向量 \([m]\) 送到 \(m\)，得到满射 \(p_0:P_0\to K_0\)。归纳地令 \(K_{n+1}=\ker p_n\)，取以 \(K_{n+1}\) 的底集合为基的自由模 \(P_{n+1}\) 及满射 \(p_{n+1}:P_{n+1}\to K_{n+1}\)。这里 \(K_1\) 是原生成元之间的关系模，\(K_2\) 是为 \(K_1\) 选择的生成元之间的关系模，后续各 \(K_n\) 以同样方式记录上一层的关系。

取增广 \(\varepsilon=p_0\)，对 \(n\ge1\) 定义
\[
\mathrm{d}_n:P_n\xrightarrow{p_n}K_n\hookrightarrow P_{n-1}.
\]
后一个箭头是单射，所以 \(\ker\mathrm{d}_n=\ker p_n=K_{n+1}\)；而 \(p_{n+1}\) 满射，所以 \(\operatorname{im}\mathrm{d}_{n+1}=K_{n+1}\)。因此每个正次数上像等于核，零次上也有 \(\operatorname{im}\mathrm{d}_1=K_1=\ker\varepsilon\)，增广满射，所构造的序列正合。

各项自由还保证它们投射。具体地，给定满射 \(e:X\to Y\) 及线性映射 \(f:R^{(S)}\to Y\)，按本卷的集合指标选择约定，为每个标准基向量 \(b_s\) 选择 \(x_s\in X\)，使 \(e(x_s)=f(b_s)\)。自由模的泛性质把 \(b_s\mapsto x_s\) 唯一延拓为线性映射 \(\widetilde f:R^{(S)}\to X\)；\(e\widetilde f\) 与 \(f\) 在基上相同，故 \(e\widetilde f=f\)。选择用在这一族原像上。对反环 \(R^{\mathrm{op}}\) 使用同一构造便得到右模版本。
\end{proof}

这个构造通常很大。对 \(\mathbb Z/m\mathbb Z\)，其中 \(m>1\)，较合适的消解是 \(0\to\mathbb Z\xrightarrow{\times m}\mathbb Z\to\mathbb Z/m\mathbb Z\to0\)。设 \(k\) 是域。对 \(R=k[\epsilon]/(\epsilon^2)\) 上的 \(k=R/(\epsilon)\)，则有无限消解
\[
\cdots\xrightarrow{\epsilon}R\xrightarrow{\epsilon}R
\xrightarrow{\epsilon}R\longrightarrow k\longrightarrow0.
\]
每个元素唯一写成 \(a+b\epsilon\)，其中 \(a,b\in k\)，并且 \(\epsilon(a+b\epsilon)=a\epsilon\)。因此乘 \(\epsilon\) 的核是所有 \(b\epsilon\)，像也是所有 \(a\epsilon\)，二者均为 \((\epsilon)\)；增广的核同样为 \((\epsilon)\)，所以该序列正合。这个有限生成模的消解仍有无限多个非零项，有限生成本身不使逐层取关系的过程停下来。

\subsection{内射消解}
\label{b4:subsec:0701:02}

投射性约束映射的出发对象：从它出发的映射能够沿满态射提升。内射性则约束映射的到达对象：映入它的映射能够沿单态射延拓。因而内射消解从嵌入原对象开始，逐层处理余核。

\begin{definition}\label{b4:def:injective-resolution}
设 \(\mathcal A\) 是交换范畴，\(I\) 为其对象。若对每个单态射 \(e:X\to Y\) 及态射 \(f:X\to I\)，都存在态射 \(\widetilde f:Y\to I\)，使 \(\widetilde f e=f\)，则称 \(I\) 为\term[neisheduixiang]{内射对象}[injective object]。给定对象 \(N\in\mathcal A\) 及增广上链复形
\[
0\longrightarrow N\xrightarrow{\eta}I^0\xrightarrow{\mathrm{d}^0}I^1
\xrightarrow{\mathrm{d}^1}I^2\longrightarrow\cdots
\]
若此序列正合且各 \(I^n\)（\(n\ge0\)）内射，则称它为 \(N\) 的\term[neishexiaojie]{内射消解}[injective resolution]。
\end{definition}

这里的增广项 \(N\) 不属于 \(I^\bullet\)；后者位于非负上链次数，增广 \(\eta\) 诱导 \(N\cong\ker\mathrm{d}^0=H^0(I^\bullet)\)，正次数上同调为零。构造时，先选到内射对象的单态射 \(j^0:N\hookrightarrow I^0\)，令 \(C^0=\operatorname{coker}j^0\)，再选到内射对象的单态射 \(j^1:C^0\hookrightarrow I^1\)。其余核记为 \(C^1\)，如此继续，并令
\[
\mathrm{d}^n:I^n\xrightarrow{q^n}C^n\xrightarrow{j^{n+1}}I^{n+1},
\qquad C^n=\operatorname{coker}j^n,
\]
其中 \(n\ge0\)，\(q^n\) 是余核映射，各 \(I^n\) 都内射。因为 \(j^{n+1}\) 单，\(\ker\mathrm{d}^n=\ker q^n=\operatorname{im}j^n\)；对 \(n\ge1\)，因为前一步的 \(q^{n-1}\) 满，\(\operatorname{im}\mathrm{d}^{n-1}=\operatorname{im}j^n\)。零次的核是 \(\operatorname{im}j^0\)，所以这条增广序列正合。

\ProseParagraphBreak
% 跨卷引用目标：b2:thm:enough-injectives（7.2.8）。
对任意含幺环 \(R\)，第二卷定理7.2.8已经证明每个幺性左 \(R\) 模都能嵌入内射模。这个构造先用内射交换群 \(\mathbb Q/\mathbb Z\) 分离非零元素：对 \(0\ne n\in N\)，总有加法群同态 \(\chi:N\to\mathbb Q/\mathbb Z\) 使 \(\chi(n)\ne0\)。再通过余诱导把交换群中的延拓变为 \(R\) 模中的延拓，得到内射左模
\[
J=\operatorname{Hom}_{\mathbb Z}(R,\mathbb Q/\mathbb Z),
\qquad (rf)(s)=f(sr).
\]
每个 \(\chi\) 对应左 \(R\) 线性映射 \(n\mapsto(s\mapsto\chi(sn))\)。单个映射未必分离所有元素，把全部这样的映射组成积映射，便得到嵌入
\[
N\longrightarrow\prod_{\chi\in\operatorname{Hom}_{\mathbb Z}(N,\mathbb Q/\mathbb Z)}J,
\qquad n\longmapsto\bigl(s\longmapsto\chi(sn)\bigr)_\chi.
\]
若 \(n\ne0\)，取分离它的 \(\chi\)，在 \(s=1\) 处这一分量非零，所以积映射单射。各 \(J\) 的内射性和内射模积的内射性都已在第二卷证明。本章把这个结果反复用于余核模，就得到内射消解；幺性右模的情形对反环应用同一结果。

\begin{example}\label{b4:exm:integer-injective-resolution}
作为 \(\mathbb Z\) 模，\(\mathbb Z\) 有内射消解
\[
0\longrightarrow\mathbb Z\longrightarrow\mathbb Q
\longrightarrow\mathbb Q/\mathbb Z\longrightarrow0.
\]
% 跨卷引用目标：b2:thm:divisible-injective（7.2.5）。
正合性由商群定义给出。\(\mathbb Q\) 中每个元素都能除以任意非零整数；在 \(\mathbb Q/\mathbb Z\) 中，对 \(q+\mathbb Z\) 和整数 \(m\ne0\)，有 \(m(q/m+\mathbb Z)=q+\mathbb Z\)。所以两项都可除，由第二卷定理7.2.5的可除群判据，它们都是内射模。

它与 \(\mathbb Z\) 自身的长度零投射消解适用于 Hom 的不同变量：对投射模 \(P\)，\(\operatorname{Hom}_{\mathbb Z}(P,-)\) 正合；对内射模 \(I\)，\(\operatorname{Hom}_{\mathbb Z}(-,I)\) 正合。对一般含幺环 \(R\) 及幺性左模 \(M,N\)，\secref{b4:sec:0802}将用第一变量的投射消解或第二变量的内射消解计算 \(\operatorname{Ext}_R^n(M,N)\)，并证明两种计算相同。
\end{example}

\subsection{足够多投射或内射对象的假设}
\label{b4:subsec:0701:03}

\begin{definition}\label{b4:def:enough-projectives-injectives}
设 \(\mathcal A\) 是交换范畴。若对每个对象 \(M\) 都存在投射对象 \(P\) 及满态射 \(P\twoheadrightarrow M\)，就说 \(\mathcal A\) 有\term[zugouduo-tousheduixiang]{足够多投射对象}[enough projectives]；若对每个对象 \(M\) 都存在内射对象 \(I\) 及单态射 \(M\hookrightarrow I\)，就说它有\term[zugouduo-neisheduixiang]{足够多内射对象}[enough injectives]。
\end{definition}

在前一条件下，对核反复选择投射满态射，就得到投射消解；在后一条件下，对余核反复选择内射单态射，就得到内射消解。这是\cref{b4:prop:free-resolution-exists}的范畴形式，正合性仍由每一步的核或余核保证。为每个对象选定一条消解，并不会自动得到函子：对象间的态射还需要提升到所选消解之间，提升也未必唯一。下一节的比较定理将证明提升存在且在链同伦意义下唯一；施加加性函子后再取同调，链同伦的提升给出相同的态射。这正是后续同调不变量能够成为函子的原因。

\ProseParagraphBreak
% 跨卷引用目标：b1:thm:fga-structure（9.3.1）。
交换范畴本身的公理没有包括这两项条件。例如有限交换群组成交换范畴，却没有非零投射对象。事实上，按第一卷定理9.3.1的有限生成交换群结构定理，非零有限交换群可分解为
\[
P\cong\bigoplus_i\mathbb Z/p_i^{a_i}\mathbb Z,
\qquad a_i\ge1,
\]
其中 \(p_i\) 为素数。假如 \(P\) 投射，任取一个非零直和因子 \(T=\mathbb Z/p^a\mathbb Z\)。对从 \(T\) 出发的映射，先与 \(P\to T\) 的投影复合，用 \(P\) 的投射性提升，再限制到 \(T\)，故 \(T\) 也投射。于是自然满射
\[
\mathbb Z/p^{a+1}\mathbb Z\longrightarrow\mathbb Z/p^a\mathbb Z
\]
应有截面 \(s\)。但 \(s(\bar1)\) 必须是 \(\bar1\) 的原像，可写成 \(1+p^at\pmod{p^{a+1}}\)；它不被 \(p\) 整除，所以阶为 \(p^{a+1}\)。另一方面，同态性要求
\[
p^a s(\bar1)=s(p^a\bar1)=0,
\]
所以它的阶又必须整除 \(p^a\)，矛盾。因此有限交换群范畴没有非零投射对象，不能把模范畴中的存在性结论无条件搬到它的交换子范畴。

% !TeX root = ../../Book4.tex
% 纲要标题：### 8.2 Ext
% 纲要位置：计划纲要.md，第 2656 行。

\section{Ext}
\label{b4:sec:0802}

Hom 的两个变量具有相反的箭头方向，分别用投射或内射消解得到的计算必须比较。平衡性使两种模型计算同一组 Ext，并保留两变量的连接态射。


% 纲要内容（仅作写作计划，尚非教材正文）：
%
% * Hom 的右导出函子
% * 反变与协变变量
% * Ext 长正合列及具体计算
%

\subsection{Hom 的右导出函子}
\label{b4:subsec:0802:01}

\begin{definition}\label{b4:def:ext}
设 \(R\) 是含幺环，\(M,N\) 是左 \(R\) 模。对 \(N\) 取内射消解 \(N\to I^\bullet\)，定义交换群
\[
\operatorname{Ext}_R^n(M,N)\coloneqq
H^n\bigl(\operatorname{Hom}_R(M,I^\bullet)\bigr),\qquad n\ge0.
\]
这就是左正合函子 \(\operatorname{Hom}_R(M,-)\) 的右导出函子。
\end{definition}
\symidx{\(\operatorname{Ext}_R^n(M,N)\)：Ext 群}

零次是 \(\operatorname{Hom}_R(M,N)\)。在一般非交换环上，Ext 首先是交换群，不自动是左 \(R\) 模；若 \(R\) 交换，标量乘法与所有微分相容，便得到自然的 \(R\) 模结构。下标的环也不能略去，例如 \(\operatorname{Ext}_{\mathbb Z}^1(\mathbb Z/2,\mathbb Z/2)\) 与 \(\operatorname{Ext}_{\mathbb F_2}^1(\mathbb F_2,\mathbb F_2)\) 不同。

\ProseParagraphBreak
固定 \(N\) 时，\(\operatorname{Hom}_R(-,N)\) 是从左模反范畴到交换群的左正合函子。左模的投射对象正是反范畴的内射对象，所以投射消解也能给出一套右导出函子。两套计算是否一致，需要证明，不能由名字相同得出。

\subsection{反变与协变变量}
\label{b4:subsec:0802:02}

\begin{theorem}[Ext 的平衡性]\label{b4:thm:ext-balanced}
设 \(R\) 是含幺环，\(M,N\) 是幺性左 \(R\) 模，\(P_\bullet\to M\) 为投射消解，\(N\to I^\bullet\) 为内射消解。对整数 \(n\ge0\) 取
\[
\operatorname{Hom}_R(P_\bullet,N)^n=\operatorname{Hom}_R(P_n,N),
\qquad \delta^n(f)=f\,\mathrm{d}_{P,n+1}.
\]
则有自然同构
\[
H^n\bigl(\operatorname{Hom}_R(P_\bullet,N)\bigr)
\cong\operatorname{Ext}_R^n(M,N).
\]
它对 \(M\) 反变，对 \(N\) 协变，并与两变量的连接态射相容。
\end{theorem}
\begin{proof}
令 \(D^{p,q}=\operatorname{Hom}_R(P_p,I^q)\)，其中 \(p,q\ge0\)，定义
\[
\mathrm{d}_h(f)=f \mathrm{d}_P,\qquad \mathrm{d}_v(f)=(-1)^p \mathrm{d}_I f.
\]
两次水平或竖直微分为零；一水平一竖直的两种复合相差一个负号，故 \(\mathrm{d}_h \mathrm{d}_v+\mathrm{d}_v\mathrm{d}_h=0\)。总次数为 \(p+q\)，每条对角线有限，因而总复形只用有限直和。

固定 \(p\)，\(P_p\) 的投射性使 \(\operatorname{Hom}_R(P_p,-)\) 正合，所以竖列只在 \(q=0\) 有上同调 \(\operatorname{Hom}_R(P_p,N)\)。固定 \(q\)，\(I^q\) 的内射性使 \(\operatorname{Hom}_R(-,I^q)\) 正合，所以横行只在 \(p=0\) 有上同调 \(\operatorname{Hom}_R(M,I^q)\)。增广分别给出两条复形映射
\[
\operatorname{Hom}_R(P_\bullet,N)\longrightarrow\operatorname{Tot}D,
\qquad
\operatorname{Hom}_R(M,I^\bullet)\longrightarrow\operatorname{Tot}D.
\]
在第一条的 \(q=0\) 边，竖直微分为零；在第二条的 \(p=0\) 边，水平微分为零。因此它们确实与总微分相容。\cref{b4:cor:double-complex-comparison}分别沿竖列和横行说明这两条映射都是拟同构。该结果通过逐个消去有限对角线上的分量证明。

投射、内射比较映射使上述两条映射对 \(M,N\) 自然；不同选择同伦，故同调态射确定。对短正合列选择马蹄消解后，所有总复形也组成逐次短正合列；两侧增广给出它们之间的交换图。\cref{b4:thm:cohomology-long-exact}的自然性于是保证同构与连接态射相容。
\end{proof}

为了方便表示上同调类，本定理使用无符号的 \(\delta(f)=f\mathrm{d}\)。若把 \(P_n\) 放在上链次数 \(-n\)，按一般内部 Hom 规则会得到 \(\delta(f)=(-1)^{n+1}f\mathrm{d}\)；每个次数 \(n\) 乘 \((-1)^{n(n+1)/2}\) 给出这两种复形的同构。本节以后用前一种计算约定。

\ProseParagraphBreak
给定 \(u:M'\to M\)，其投射消解提升经预合成给出 \(u^*:\operatorname{Ext}_R^n(M,N)\to\operatorname{Ext}_R^n(M',N)\)；给定 \(v:N\to N'\)，后合成给出 \(v_*\)。预合成与后合成交换，所以 Ext 是两个变量的函子，而非仅对每个固定模分别定义的一组交换群。

\subsection{Ext 长正合列及具体计算}
\label{b4:subsec:0802:03}

\begin{theorem}\label{b4:thm:ext-long-exact}
设 \(R\) 是含幺环。左 \(R\) 模短正合列 \(0\to A\to B\to C\to0\) 对任意左模 \(N\) 给出自然长正合列
\[
\begin{aligned}
0&\longrightarrow\operatorname{Hom}_R(C,N)\longrightarrow\operatorname{Hom}_R(B,N)
\longrightarrow\operatorname{Hom}_R(A,N)\\
&\longrightarrow\operatorname{Ext}_R^1(C,N)\longrightarrow\operatorname{Ext}_R^1(B,N)
\longrightarrow\operatorname{Ext}_R^1(A,N)\longrightarrow\cdots.
\end{aligned}
\]
对任意左模 \(M\)，左模短正合列 \(0\to N'\to N\to N''\to0\) 给出自然长正合列
\[
\begin{aligned}
0&\longrightarrow\operatorname{Hom}_R(M,N')\longrightarrow\operatorname{Hom}_R(M,N)
\longrightarrow\operatorname{Hom}_R(M,N'')\\
&\longrightarrow\operatorname{Ext}_R^1(M,N')\longrightarrow\operatorname{Ext}_R^1(M,N)
\longrightarrow\operatorname{Ext}_R^1(M,N'')\longrightarrow\cdots.
\end{aligned}
\]
\end{theorem}
\begin{proof}
第一列取 \(N\) 的内射消解，逐次施加反变正合函子 \(\operatorname{Hom}_R(-,I^q)\)，得到 \(0\to\operatorname{Hom}(C,I)\to\operatorname{Hom}(B,I)\to\operatorname{Hom}(A,I)\to0\)，再取同调长正合列。第二列取 \(M\) 的投射消解，对给定序列逐次施加正合函子 \(\operatorname{Hom}_R(P_p,-)\)，取长正合列，再用\cref{b4:thm:ext-balanced}识别各项。增广在零次给出通常 Hom 映射，连接态射与选择独立且自然。
\end{proof}

\begin{example}\label{b4:exm:ext-cyclic}
设 \(R\) 是主理想整环，\(0\ne a\in R\)，\(N\) 是 \(R\) 模。把 \(R/(a)\) 的两项消解送入 Hom，得到 \(N\xrightarrow{a}N\)，故
\[
\operatorname{Ext}_R^0(R/(a),N)=\{\,n\in N\mid an=0\,\},
\qquad\operatorname{Ext}_R^1(R/(a),N)=N/aN,
\]
二次及以上为零。特别地，\(\operatorname{Ext}_{\mathbb Z}^1(\mathbb Z/m,\mathbb Z/n)\cong\mathbb Z/\gcd(m,n)\mathbb Z\)，其中 \(m,n>0\)。若改取 \(N=\mathbb Q\)，乘 \(m\) 为同构，零次及正次均为零；这与 \(\mathbb Q\) 内射一致。
\end{example}

在 \(R=k[\epsilon]/(\epsilon^2)\) 上，对\secref{b4:sec:0701}给出的 \(k\) 的周期消解施加 \(\operatorname{Hom}_R(-,k)\)，所有微分为零。因此 \(\operatorname{Ext}_R^n(k,k)\cong k\) 对每个 \(n\ge0\) 成立。相同的底层向量空间，在底环为 \(k\) 时正次 Ext 全为零；记录底环是计算的一部分。

% !TeX root = ../../Book4.tex
% 纲要标题：### 8.3 马蹄引理
% 纲要标签：b4:sec:0703
% 纲要位置：计划纲要.md，第 3287 行。

\section{马蹄引理}
\label{b4:sec:0703}

短正合列的两端已有消解时，中项的消解仍须保存原扩张的不分裂信息。马蹄构造把这部分信息放进微分的非对角项，使三份消解逐次相容。


% 纲要内容（仅作写作计划，尚非教材正文）：
%
% * 短正合列与相容消解
% * 投射及内射版本
% * 消解上的长正合列
%

\subsection{短正合列与相容消解}
\label{b4:subsec:0703:01}

假设已有短正合列 \(0\to A\xrightarrow{i}B\xrightarrow{q}C\to0\)，又已分别选好 \(A,C\) 的投射消解。希望构造 \(B\) 的消解时，最自然的逐次对象是 \(P_n^A\oplus P_n^C\)。困难在于微分通常不能取两个微分的直和：那样得到的零次同调是 \(A\oplus C\)，未必是 \(B\)。原列的不分裂信息必须进入微分的非对角项。

\begin{lemma}\label{b4:lem:horseshoe-step}
设 \(\mathcal A\) 是交换范畴，\(0\to A\xrightarrow{i}B\xrightarrow{q}C\to0\) 正合，\(p_A:P_A\to A\)、\(p_C:P_C\to C\) 满，且 \(P_C\) 投射。选 \(t:P_C\to B\) 使 \(qt=p_C\)，令 \(p_B=(ip_A,t):P_A\oplus P_C\to B\)。则 \(p_B\) 满，且
\[
0\longrightarrow\ker p_A\longrightarrow\ker p_B
\longrightarrow\ker p_C\longrightarrow0
\]
正合。
\end{lemma}
\begin{proof}
在两行分别为 \(0\to P_A\to P_A\oplus P_C\to P_C\to0\) 和给定短正合列的图中，竖直态射取 \(p_A,p_B,p_C\)。图由 \(qt=p_C\)、\(qi=0\) 交换：
\[
\begin{tikzcd}[column sep=2.5em,row sep=large]
0\arrow[r]&P_A\arrow[r]\arrow[d,two heads,"p_A"']&P_A\oplus P_C\arrow[r]\arrow[d,"p_B"]&P_C\arrow[r]\arrow[d,two heads,"p_C"]&0\\
0\arrow[r]&A\arrow[r,"i"']&B\arrow[r,"q"']&C\arrow[r]&0
\end{tikzcd}
\]
由\cref{b4:thm:snake-lemma}得到正合列
\[
\begin{aligned}
0\longrightarrow\ker p_A\longrightarrow\ker p_B
&\longrightarrow\ker p_C\longrightarrow\operatorname{coker}p_A\\
&\longrightarrow\operatorname{coker}p_B
\longrightarrow\operatorname{coker}p_C\longrightarrow0.
\end{aligned}
\]
因 \(p_A,p_C\) 均满，两端余核为零。于是 \(\ker p_B\to\ker p_C\) 满，且 \(\operatorname{coker}p_B=0\)，这正是结论。
\end{proof}

\subsection{马蹄引理的投射与内射版本}
\label{b4:subsec:0703:02}

\begin{theorem}[马蹄引理]\label{b4:thm:horseshoe}
设 \(\mathcal A\) 是交换范畴，\(0\to A\to B\to C\to0\) 是短正合列。给定 \(A,C\) 的投射消解 \(P^A_\bullet,P^C_\bullet\)，存在 \(B\) 的投射消解 \(P^B_\bullet\)，使增广相容，且
\[
0\longrightarrow P^A_\bullet\longrightarrow P^B_\bullet
\longrightarrow P^C_\bullet\longrightarrow0
\]
逐次分裂正合，并可取 \(P_n^B=P_n^A\oplus P_n^C\)。对偶地，给定 \(A,C\) 的内射消解，可取 \(I_B^n=I_A^n\oplus I_C^n\)，得到相容且逐次分裂的短正合列 \(0\to I_A^\bullet\to I_B^\bullet\to I_C^\bullet\to0\)。此结果称为\term[matiyinli]{马蹄引理}[horseshoe lemma]。
\end{theorem}
\begin{proof}
这里的分裂逐次数成立，截面未必与微分交换。选定双积表示后，中间微分的非对角块正是需要另外构造的部分；只有存在复形截面时，才能用它选取双积表示，使这些非对角块同时为零。

先将\cref{b4:lem:horseshoe-step}用于 \(P_0^A\to A\)、\(P_0^C\to C\)，得到 \(P_0^B\to B\) 及其核列。记这三个核为 \(K_1^A,K_1^B,K_1^C\)。给定消解中的下一项分别满射到 \(K_1^A,K_1^C\)，于是再用引理得到 \(P_1^B\to K_1^B\)，其核又组成短正合列。归纳继续。把 \(P_n^B\to K_n^B\) 与 \(K_n^B\hookrightarrow P_{n-1}^B\) 复合，就是中间消解的微分；每一步定义保证像等于下一个核。

各次包含与投影通过核图构造，因而与微分交换。其对象列是双积的分裂列；\(P_n^A\oplus P_n^C\) 投射，因为两个分量的态射可以分别提升后相加。这样得到全部结论。内射版本在 \(\mathcal A^{\mathrm{op}}\) 应用已证投射版本，并反转箭头与链指标；有限双积在反范畴仍是双积，故逐次分裂性不变。
\end{proof}

选定逐次双积表示以后，相容消解可以排列如下：
\[
\begin{tikzcd}[column sep=large,row sep=large]
P_1^A\arrow[r]\arrow[d,"\mathrm{d}_1^A"']&P_1^B\arrow[r]\arrow[d,"\mathrm{d}_1^B"]&P_1^C\arrow[d,"\mathrm{d}_1^C"]\\
P_0^A\arrow[r]\arrow[d]&P_0^B\arrow[r]\arrow[d]&P_0^C\arrow[d]\\
A\arrow[r]&B\arrow[r]&C
\end{tikzcd}
\]
从 \(P_0^A,P_0^B,P_0^C\) 向下的箭头为增广；上面两排来自逐次分裂的消解短正合列。微分具体为
\[
\mathrm{d}_n^B=\begin{pmatrix}\mathrm{d}_n^A&h_n\\0&\mathrm{d}_n^C\end{pmatrix},
\qquad \mathrm{d}_{n-1}^Ah_n+h_{n-1}\mathrm{d}_n^C=0.
\]
第二个等式在 \(n\ge2\) 时成立，是 \((\mathrm{d}^B)^2=0\) 的非对角分量，显示原来的扩张如何影响消解。零次的相容关系则由增广给出。

\subsection{消解上的长正合列}
\label{b4:subsec:0703:03}

\begin{proposition}\label{b4:prop:horseshoe-functor-sequence}
设 \(\mathcal A,\mathcal B\) 为交换范畴，\(F:\mathcal A\to\mathcal B\) 为加性函子，\(0\to A\to B\to C\to0\) 为 \(\mathcal A\) 中的短正合列。给定两端 \(A,C\) 的投射消解，并按马蹄引理构造相容的消解短正合列 \(0\to P^A_\bullet\to P^B_\bullet\to P^C_\bullet\to0\)。此列经 \(F\) 逐次作用后仍是逐次分裂的复形短正合列，因此其同调对象有自然连接态射及长正合列。两端取内射消解及相应的马蹄构造时，同样成立。
\end{proposition}
\begin{proof}
分裂短正合列由态射等式 \(ri=1\)、\(ps=1\)、\(ir+sp=1\) 表示。加性函子保持这些等式，故不必假定 \(F\) 正合。再用\cref{b4:thm:cohomology-long-exact}，在链情形反转指标，即得到长正合列。

为证明自然性，给定短正合列间的交换图
\[
\begin{tikzcd}[column sep=large,row sep=large]
0\arrow[r]&A\arrow[r,"i"]\arrow[d,"a"']
&B\arrow[r,"q"]\arrow[d,"b"]
&C\arrow[r]\arrow[d,"c"]&0\\
0\arrow[r]&A'\arrow[r,"i'"']
&B'\arrow[r,"q'"']&C'\arrow[r]&0.
\end{tikzcd}
\]
源端的马蹄消解记为 \(P^A,P^B,P^C\)，目标端记为 \(Q^A,Q^B,Q^C\)。选定分别提升 \(a,c\) 的端点比较映射 \(u:P^A\to Q^A\)、\(v:P^C\to Q^C\)。用双积表示两套微分，分别记非对角项为 \(h_n\)、\(h'_n\)。设各端点增广为 \(\varepsilon_{P^A},\varepsilon_{P^C},\varepsilon_{Q^A},\varepsilon_{Q^C}\)，中间增广写成
\[
\begin{aligned}
\varepsilon_P&=(i\varepsilon_{P^A},t_P):
P_0^A\oplus P_0^C\longrightarrow B,\\
\varepsilon_Q&=(i'\varepsilon_{Q^A},t_Q):
Q_0^A\oplus Q_0^C\longrightarrow B',
\end{aligned}
\]
其中 \(qt_P=\varepsilon_{P^C}\)、\(q't_Q=\varepsilon_{Q^C}\)。我们构造中间比较映射
\[
w_n=\begin{pmatrix}u_n&y_n\\0&v_n\end{pmatrix}.
\]
零次需要满足 \(\varepsilon_Qw_0=b\varepsilon_P\)。第一列已由 \(u_0\) 的增广条件满足；第二列要求
\[
i'\varepsilon_{Q^A}y_0=bt_P-t_Qv_0:
P_0^C\longrightarrow B'.
\]
右端与 \(q'\) 的复合为
\[
q'(bt_P-t_Qv_0)
=c\varepsilon_{P^C}-\varepsilon_{Q^C}v_0=0.
\]
因 \(i'\) 是 \(q'\) 的核，右端唯一分解为 \(i'z_0\)，其中 \(z_0:P_0^C\to A'\)。由 \(P_0^C\) 投射，沿满态射 \(\varepsilon_{Q^A}:Q_0^A\to A'\) 提升 \(z_0\)，得到 \(y_0:P_0^C\to Q_0^A\)。这就满足了中间的增广条件。

已构造到次数 \(n-1\) 后，次数 \(n\) 的链映射条件只剩非对角块
\[
\mathrm{d}_n^{Q^A}y_n=r_n,
\qquad r_n=u_{n-1}h_n+y_{n-1}\mathrm{d}_n^{P^C}-h'_nv_n.
\]
这里 \(r_n:P_n^C\to Q_{n-1}^A\)，其余三块由 \(u,v\) 是链映射及零块自动满足。当 \(n=1\) 时，源、目标增广与微分的相容关系分别给出
\[
i\varepsilon_{P^A}h_1+t_P\mathrm{d}_1^{P^C}=0,
\qquad
i'\varepsilon_{Q^A}h'_1+t_Q\mathrm{d}_1^{Q^C}=0.
\]
结合已构造的 \(y_0\)，可得
\[
\begin{aligned}
i'\varepsilon_{Q^A}r_1
&=-bt_P\mathrm{d}_1^{P^C}
 +(bt_P-t_Qv_0)\mathrm{d}_1^{P^C}
 +t_Q\mathrm{d}_1^{Q^C}v_1\\
&=0,
\end{aligned}
\]
最后一步使用 \(v_0\mathrm{d}_1^{P^C}=\mathrm{d}_1^{Q^C}v_1\)。因 \(i'\) 单，\(\varepsilon_{Q^A}r_1=0\)。当 \(n\ge2\) 时，利用两套微分的非对角关系和上一次数的链映射等式，得到
\[
\begin{aligned}
\mathrm{d}_{n-1}^{Q^A}r_n
={}&-u_{n-2}h_{n-1}\mathrm{d}_n^{P^C}\\
&+\bigl(u_{n-2}h_{n-1}
 +y_{n-2}\mathrm{d}_{n-1}^{P^C}
 -h'_{n-1}v_{n-1}\bigr)\mathrm{d}_n^{P^C}\\
&+h'_{n-1}\mathrm{d}_n^{Q^C}v_n
=0.
\end{aligned}
\]
其中含 \(y_{n-2}\) 的项由 \((\mathrm{d}^{P^C})^2=0\) 消失，其余项由 \(v\) 的链映射等式抵消。因此 \(r_n\) 分解到目标核；目标消解正合，使 \(Q_n^A\) 满射到该核，再由 \(P_n^C\) 投射取提升 \(y_n\)。归纳给出消解短正合列之间的严格交换图。

对图施加 \(F\)，再用复形长正合列的自然性，即有
\[
H_{n-1}(F(u))\,\delta_P
=\delta_Q\,H_n(F(v)).
\]
若固定这两套消解而改变比较映射或中间提升，则所得三个比较链映射仍分别提升 \(a,b,c\)。由\cref{b4:prop:comparison-homotopy-unique}，它们分别链同伦；加性函子保持这些同伦，故诱导的同调态射不变。

若改变端点消解或马蹄消解，则把上述构造用于原短正合列的恒等图。三个比较映射都提升恒等，由\cref{b4:cor:resolutions-homotopy-equivalent}在施加 \(F\) 后诱导同调同构。比较映射的同伦唯一性保证这些同构不依赖提升，并满足传递律；刚才的自然性等式保证它们也与连接态射交换。因此，不同消解所得的长正合列在这些典范比较同构下相容，而不是在字面上具有相同的同调对象。内射情形在反范畴作同一构造，转回后得到上同调长正合列及其相容性。
\end{proof}

连接态射有可操作的含义：把商复形中的圈提升到中间复形，对提升取微分，再把所得像读成左端复形中的圈。马蹄引理保证即使函子不正合，仍有逐次正合的复形列来执行这一过程；下一章的 Ext 与 Tor 正是由此获得正合列。

% !TeX root = ../../Book4.tex
% 纲要标题：### 8.4 特殊消解
% 纲要标签：b4:sec:0704
% 纲要位置：计划纲要.md，第 3392 行。

\section{特殊消解}
\label{b4:sec:0704}

消解的长度与形式会随底环改变：主理想整环给出短自由消解，正则列给出 Koszul 消解，群或代数乘法给出 bar 消解。每种公式都需要独立检查其正合性。


% 纲要内容（仅作写作计划，尚非教材正文）：
%
% * 主理想整环模的短消解
% * Koszul 消解简介
% * 群代数和代数的 bar 消解预告
%
% ---
%

\subsection{主理想整环模的短消解}
\label{b4:subsec:0704:01}

% 跨卷引用目标：b2:thm:pid-free-submodule（4.3.1）。
第二卷定理4.3.1证明了主理想整环上有限秩自由模的子模自由。有限秩证明逐个加入基向量，每一步用一个主理想描述新增加的部分。无限基未必能按自然数排成一列，也没有供有限归纳结束的最后一项。这里使用选择公理的良序形式，把任意基按序数排列：有直接前驱的阶段只增加一个基向量，没有直接前驱的极限阶段则取此前各阶段的并。自由模中的向量只有有限个非零坐标，因而取并仍只使用有限线性组合。

\begin{theorem}\label{b4:thm:pid-free-submodule}
设 \(R\) 是主理想整环，\(F\) 是任意秩自由 \(R\) 模。每个子模 \(N\subseteq F\) 都自由。因此每个 \(R\) 模 \(M\) 都有长度至多为一的自由消解 \(0\to F_1\to F_0\to M\to0\)。
\end{theorem}
\begin{proof}
用选择公理良序排列 \(F\) 的一组基，写成 \((e_\alpha)_{\alpha<\kappa}\)。对 \(\beta\le\kappa\)，令 \(F_\beta\) 由 \(\alpha<\beta\) 的基向量生成，\(N_\beta=N\cap F_\beta\)。在增加 \(e_\alpha\) 的阶段，取其坐标映射
\[
\pi_\alpha:N_{\alpha+1}\longrightarrow R.
\]
其核恰为 \(N_\alpha\)，像 \(I_\alpha\) 是 \(R\) 的一个理想，故模的第一同构定理给出
\[
N_{\alpha+1}/N_\alpha\cong I_\alpha.
\]
若 \(I_\alpha\ne0\)，取非零生成元 \(a_\alpha\)，并在 \(N_{\alpha+1}\) 选一个 \(e_\alpha\) 坐标为 \(a_\alpha\) 的元素 \(x_\alpha\)。任意 \(y\in N_{\alpha+1}\) 的该坐标是 \(ra_\alpha\)，所以 \(y-rx_\alpha\in N_\alpha\)，这证明 \(N_{\alpha+1}=N_\alpha+Rx_\alpha\)。若 \(rx_\alpha\in N_\alpha\)，看 \(e_\alpha\) 坐标得 \(ra_\alpha=0\)。因为 \(R\) 是整环且 \(a_\alpha\ne0\)，必有 \(r=0\)，所以两部分交为零，即
\[
N_{\alpha+1}=N_\alpha\oplus Rx_\alpha.
\]
若 \(I_\alpha=0\)，则 \(N_{\alpha+1}=N_\alpha\)。

在非零极限序数 \(\lambda\) 处，给定非零 \(y\in N_\lambda\)，它的非零坐标指标构成有限集，因而有最大元 \(\gamma<\lambda\)。极限序数没有直接前驱，所以 \(\gamma+1<\lambda\)，且 \(y\in N_{\gamma+1}\)。零向量也属于此前各阶段，故
\[
N_\lambda=\bigcup_{\beta<\lambda}N_\beta.
\]
从 \(N_0=0\) 开始，上述后继阶段的分解与极限阶段的取并表明，每个 \(N_\beta\) 都由已选出的 \(x_\alpha\)（\(\alpha<\beta\)）生成，特别是它们生成 \(N=N_\kappa\)。若有有限关系
\[
r_1x_{\alpha_1}+\cdots+r_mx_{\alpha_m}=0,
\qquad \alpha_1<\cdots<\alpha_m,
\]
看 \(e_{\alpha_m}\) 坐标：前面各项在该坐标为零，故 \(r_ma_{\alpha_m}=0\)，从而 \(r_m=0\)。再对余下的有限项重复，得到全部系数为零。因此这些 \(x_\alpha\) 构成 \(N\) 的基。

最后取自由满射 \(F_0\twoheadrightarrow M\)，其核 \(F_1\subseteq F_0\) 由刚证明的结论自由，于是
\[
0\longrightarrow F_1\longrightarrow F_0\longrightarrow M\longrightarrow0
\]
是所需的长度至多为一的自由消解。
\end{proof}

有限生成 \(M\) 时可先取有限秩 \(F_0\)，并由有限秩版本得到有限秩 \(F_1\)。若 \(M=R/(a)\)、\(a\ne0\)，其短消解就是乘 \(a\) 的映射；整环假设在这里保证该映射单射。

\subsection{Koszul 消解简介}
\label{b4:subsec:0704:02}

设 \(R\) 为交换含幺环，\(x\in R\)。两项链复形 \(K(x)=(0\to R\xrightarrow{x}R\to0)\) 把两个 \(R\) 分别置于次数一、零，微分 \(\mathrm{d}_1\) 是乘 \(x\)。其零次同调为余核 \(R/(x)\)，一次同调为核 \(\{\,r\in R\mid xr=0\,\}\)。因此它成为 \(R/(x)\) 的消解，准确地要求乘 \(x\) 单射。

\begin{definition}\label{b4:def:koszul-preview}
设 \(R\) 是交换含幺环，\(r\ge0\) 为整数，\(x_1,\ldots,x_r\in R\)。取自由 \(R\) 模 \(R^r\) 的基 \(e_1,\ldots,e_r\)，令 \(K_n=\bigwedge^nR^r\)（\(0\le n\le r\)），其余次数取零，并用 \(R\) 线性延拓定义微分
\[
\mathrm{d}(e_{i_1}\wedge\cdots\wedge e_{i_n})
=\sum_{j=1}^n(-1)^{j-1}x_{i_j}
 e_{i_1}\wedge\cdots\widehat{e_{i_j}}\cdots\wedge e_{i_n}.
\]
所得链复形记为 \(K_\bullet(x_1,\ldots,x_r;R)\)，称为\term[koszul-fuxing]{Koszul 复形}[Koszul complex]。
\end{definition}
\symidx{\(K_\bullet(x_1,\ldots,x_r;R)\)：Koszul 复形}

例如，当 \(i<j\) 时，
\[
\mathrm{d}(e_i\wedge e_j)=x_ie_j-x_je_i,
\qquad
\mathrm{d}^2(e_i\wedge e_j)=x_ix_j-x_jx_i=0.
\]
% 跨卷引用目标：b3:thm:koszul-regular-resolution（15.3.2）。
一般情形也如此：删去两个基向量的两种次序符号相反，而 \(R\) 的交换性使两种次序中的系数相同，所以 \(\mathrm{d}^2=0\)。外幂及楔积的定义与基已在第二卷第九章建立。第三卷定理15.3.2证明了正则序列给出 Koszul 消解；下面利用同调长正合列给出第二证明，考察每增加一个元素时同调如何变化。

\begin{proposition}\label{b4:prop:koszul-resolution}
% 跨卷引用目标：b3:def:regular-sequence（15.1.1）。
设 \(R\) 是交换含幺环，\(r\ge0\) 为整数，\(x_1,\ldots,x_r\in R\)。假设 \((x_1,\ldots,x_r)\ne R\)，且对每个 \(1\le i\le r\)，乘 \(x_i\) 在 \(R/(x_1,\ldots,x_{i-1})\) 上单射，即按第三卷定义15.1.1，\(x_1,\ldots,x_r\) 是一个 \(R\) 正则序列。则 \(R\) 上的 Koszul 复形 \(K_\bullet(x_1,\ldots,x_r;R)\) 连同零次增广 \(R\to R/(x_1,\ldots,x_r)\) 是 \(R/(x_1,\ldots,x_r)\) 的有限自由消解。
\end{proposition}
\begin{proof}
\(r=0\) 时取集中于零次的 \(R\)，结论成立。对正整数 \(r\) 归纳。\(r=1\) 就是前面的两项计算。设 \(K'=K_\bullet(x_1,\ldots,x_{r-1};R)\)。楔积的每个基元要么不含 \(e_r\)，要么唯一写成前 \(r-1\) 个基向量的楔积再楔上 \(e_r\)。因此
\[
K_n=K'_n\oplus(K'_{n-1}\wedge e_r).
\]
若 \(a\in K'_n\)、\(b\in K'_{n-1}\)，则在 \(b\) 的每个楔积基元之后添入的 \(e_r\) 位于第 \(n\) 个位置，删去它的符号为 \((-1)^{n-1}\)，所以
\[
\mathrm{d}(a+b\wedge e_r)=\mathrm{d}'a+(-1)^{n-1}x_rb+(\mathrm{d}'b)\wedge e_r.
\]
令 \(Q_n=K'_{n-1}\)、\(\mathrm{d}_{Q,n}=\mathrm{d}'_{n-1}\)。这里重新编号次数并保留原微分，与前章微分带负号的移位约定不同。以
\[
\iota_n(a)=a,\qquad q_n(a+b\wedge e_r)=b
\]
定义包含与投影。上述微分公式给出
\(\mathrm{d}\iota=\iota\mathrm{d}'\) 及
\(q_{n-1}\mathrm{d}=\mathrm{d}_{Q,n}q_n\)，因此它们组成短正合复形列
\[
0\longrightarrow K'_\bullet\xrightarrow{\iota}K_\bullet
\xrightarrow{q}Q_\bullet\longrightarrow0.
\]
每个次数的投影都以 \(b\mapsto b\wedge e_r\) 为截面，但这个截面与微分的差为左端中的 \((-1)^{n-1}x_rb\)，一般并不组成复形截面。这与马蹄微分的非对角项一样，把逐次分裂与复形分裂区分开来。

对 \(n\ge1\)，\(Q_n\) 中的圈 \(b\in K'_{n-1}\) 满足 \(\mathrm{d}'b=0\)。按照同调长正合列中连接同态的构造，先在 \(K_n\) 中以 \(b\wedge e_r\) 提升它，再取微分，得到
\[
\mathrm{d}(b\wedge e_r)=(-1)^{n-1}x_rb.
\]
所得元素完全落在子复形 \(K'\) 中，故在识别 \(H_n(Q)=H_{n-1}(K')\) 后，连接同态为
\[
\delta_n:H_{n-1}(K')\longrightarrow H_{n-1}(K'),
\qquad [b]\longmapsto(-1)^{n-1}[x_rb].
\]
归纳假设给出 \(H_j(K')=0\)（\(j>0\)）及 \(H_0(K')=R/(x_1,\ldots,x_{r-1})\)。当 \(n\ge2\) 时，长正合列中的片段为
\[
H_n(K')\longrightarrow H_n(K)\longrightarrow H_n(Q)=H_{n-1}(K'),
\]
两端为零，所以 \(H_n(K)=0\)。只有 \(n=1\) 时右端仍为 \(H_0(K')\)，对应的片段为
\[
0\longrightarrow H_1(K)\longrightarrow R/(x_1,\ldots,x_{r-1})
\xrightarrow{x_r}R/(x_1,\ldots,x_{r-1})
\longrightarrow H_0(K)\longrightarrow0.
\]
正则性在这一归纳步骤中保证中间乘 \(x_r\) 的映射单射，故一次同调也消失，而零次同调为所需商环。Koszul 复形的定义、上述短正合列及连接同态公式本身都不要求正则性；这一条件控制的是正次数同调的消失。
\end{proof}

例如，设 \(k\) 为域，令 \(R=k[x,y]\)，则 \(k=R/(x,y)\) 有消解
\[
0\longrightarrow R\xrightarrow{c\mapsto(-yc,xc)}R^2
\xrightarrow{(a,b)\mapsto xa+yb}R\longrightarrow k\longrightarrow0.
\]
把此消解与 \(k\) 在 \(R\) 上张量后，\(x,y\) 在 \(k\) 上的作用均为零，因此张量后的微分全为零。下一章将由此计算出
\[
\operatorname{Tor}_1^R(k,k)\cong k^2,
\qquad \operatorname{Tor}_2^R(k,k)\cong k.
\]
这些非零同调来自张量后的复形；原来的 Koszul 消解在正次数仍然正合。这也说明 \(-\otimes_Rk\) 没有保持该消解的正合性，与\cref{b4:ex:06-tensor-breaks-quasi}中用 \(\mathbb Z/2\mathbb Z\) 张量破坏拟同构的计算一样。Tor 函子正是用这些新出现的同调记录张量函子未保持正合性的信息。

\subsection{群代数与代数的 bar 消解}
\label{b4:subsec:0704:03}

群环和一般代数的自然模往往没有短消解，却有一种由“插入一个符号”控制正合性的自由消解。这里介绍其形式；后面群上同调与 Hochschild 上同调的章节将说明它为何适合相应问题。

\begin{example}\label{b4:exm:homogeneous-bar}
设 \(G\) 是群，\(k\) 是交换含幺环。对 \(n\ge0\)，令 \(B_n\) 为以 \((g_0,\ldots,g_n)\in G^{n+1}\) 为基的自由 \(k\) 模，并让 \(G\) 同时左乘所有坐标。任意元组唯一写成
\[
(g_0,\ldots,g_n)
=g_0\cdot(1,g_0^{-1}g_1,\ldots,g_0^{-1}g_n).
\]
因此首坐标为 \(1\) 的元组构成各轨道的唯一代表，每个轨道上的自由 \(k\) 模同构于左正则模 \(kG\)，从而
\[
B_n\cong\bigoplus_{(h_1,\ldots,h_n)\in G^n}kG
\]
是自由幺性左 \(kG\) 模，其中 \(n=0\) 时右边只有一个直和项。对 \(n\ge1\) 令
\[
\mathrm{d}_n(g_0,\ldots,g_n)
=\sum_{i=0}^n(-1)^i(g_0,\ldots,\widehat{g_i},\ldots,g_n),
\qquad\varepsilon(g_0)=1.
\]
若 \(i<j\)，先删去第 \(j\) 项再删第 \(i\) 项的系数为 \((-1)^{i+j}\)，先删第 \(i\) 项再删原来的第 \(j\) 项的系数为 \((-1)^{i+j-1}\)：第二次删除时该项的编号已变为 \(j-1\)。两项相消，故 \(\mathrm{d}_{n-1}\mathrm{d}_n=0\)（\(n\ge2\)）；次数一还满足 \(\varepsilon\mathrm{d}_1(g_0,g_1)=1-1=0\)。

在底层 \(k\) 模上，用 \(k\) 线性延拓定义
\[
s_n(g_0,\ldots,g_n)=(1,g_0,\ldots,g_n)\quad(n\ge0),
\qquad s_{-1}(1)=(1).
\]
删除新插入的首项给出原元组，其余删除项与 \(s\mathrm{d}\) 抵消，因此
\[
\begin{aligned}
\mathrm{d}_{n+1}s_n+s_{n-1}\mathrm{d}_n&=1_{B_n}\quad(n\ge1),\\
\mathrm{d}_1s_0+s_{-1}\varepsilon&=1_{B_0},
\qquad \varepsilon s_{-1}=1_k.
\end{aligned}
\]
底层增广 \(k\) 模复形于是可缩，特别是正合。所有微分与增广都是 \(kG\) 线性的，核、像在遗忘 \(G\) 作用后仍是同一底层子模，故底层等式 \(\ker\mathrm{d}_n=\operatorname{im}\mathrm{d}_{n+1}\) 也就是 \(kG\) 模中的正合性；增广处同样成立。这便是平凡幺性左 \(kG\) 模 \(k\) 的自由消解。

不过，对 \(g\in G\) 一般有
\[
\begin{aligned}
s_n\bigl(g\cdot(g_0,\ldots,g_n)\bigr)&=(1,gg_0,\ldots,gg_n),\\
g\cdot s_n(g_0,\ldots,g_n)&=(g,gg_0,\ldots,gg_n).
\end{aligned}
\]
首坐标的差别说明上述 \(k\) 线性收缩通常不是 \(kG\) 线性的。正合性可由底层核、像检测，可缩性却要求收缩属于正在讨论的模范畴，因而不能由这个收缩断言增广复形在 \(kG\) 模范畴中可缩。
\end{example}

代数的 bar 构造把元组换成张量串。设 \(k\) 为域，\(A\) 为含幺 \(k\) 代数，令 \(A^e=A\otimes_k A^{\mathrm{op}}\)。对 \(n\ge0\)，取 \(B_n=A^{\otimes_k(n+2)}\)，其幺性左 \(A^e\) 模作用为
\[
\begin{aligned}
&(a\otimes b^{\mathrm{op}})\cdot(a_0\otimes\cdots\otimes a_{n+1})\\
&\hspace{2em}=aa_0\otimes a_1\otimes\cdots\otimes a_n\otimes a_{n+1}b.
\end{aligned}
\]
也就是说，首尾因子分别承担左、右 \(A\) 作用。对 \(n\ge1\)，微分交错地乘合相邻因子：
\[
\mathrm{d}_n(a_0\otimes\cdots\otimes a_{n+1})
=\sum_{i=0}^n(-1)^i
 a_0\otimes\cdots\otimes a_ia_{i+1}\otimes\cdots\otimes a_{n+1},
\]
增广 \(\varepsilon:B_0=A\otimes_k A\to A\) 为乘法。结合律使两次乘合的项成对抵消，故得到复形。在左端插入 \(1\)，即取
\[
\begin{aligned}
s_n(a_0\otimes\cdots\otimes a_{n+1})
&=1\otimes a_0\otimes\cdots\otimes a_{n+1},\\
s_{-1}(a)&=1\otimes a.
\end{aligned}
\]
便给出底层向量空间的收缩：乘合新插入的首因子产生原张量，其余项与 \(s\mathrm{d}\) 抵消。收缩只需 \(k\) 线性，而微分与增广都为 \(A^e\) 线性，故与群的情形一样，可以由底层正合性得到 \(A^e\) 模的正合性。

\ProseParagraphBreak
自由性则由首尾因子与中间因子的分离得到。取 \(A^{\otimes_k n}\) 的一组 \(k\) 基 \(\mathcal B_n\)，其中 \(A^{\otimes_k0}=k\)。映射
\[
A^e\otimes_k A^{\otimes_k n}\longrightarrow B_n,
\qquad
(a\otimes b^{\mathrm{op}})\otimes v\longmapsto a\otimes v\otimes b
\]
是左 \(A^e\) 模同构，因此
\[
B_n\cong A^e\otimes_k A^{\otimes_k n}
\cong\bigoplus_{v\in\mathcal B_n}A^e.
\]
由此得到 \(A\) 作为幺性左 \(A^e\) 模的自由消解。域的假设在此保证中间张量因子有向量空间基；在一般交换底环上不能直接由同一公式断言各项自由。双模约定、归一化及由此得到的上同调计算见\chapref{b4:ch:12}。


\begin{chaptersummary}
投射消解反复把关系写成投射对象的商，内射消解反复处理嵌入后的商对象。模范畴中的存在性使每个模都能进入这两种构造；一般交换范畴必须先检查相应假设。比较定理给出不同消解之间的同伦等价，加性函子仍保持这些等价；当目标为交换范畴时，取同调便得到典范比较同构。

\ProseParagraphBreak
马蹄消解的每次对象是双积，微分却可能含非对角项。逐次分裂性保证加性函子作用后的列仍逐次分裂；若目标为交换范畴，再取同调就得到长正合列。计算时则宜尽量选适合对象的消解：主理想整环给出两项自由消解，正则列给出有限 Koszul 消解，而 bar 构造用统一的较大模型换取态射上的便利。
\end{chaptersummary}

\BookExercises

\begin{exercise}\label{b4:ex:07-noetherian-finite-resolution}
设 \(R\) 是左 Noether 的含幺环\footnote{埃米·诺特（Amalie Emmy Noether，1882-1935），德国数学家，建立现代抽象代数中理想的链条件理论，并提出连接对称性与守恒律的定理。}，\(M\) 是有限生成幺性左 \(R\)-模。证明可以给 \(M\) 选择每次秩都有限的自由消解。该证明是否保证只有有限个非零项？
\end{exercise}
\begin{solution}
有限生成性给出 \(R^{r_0}\to M\) 的满射。因 \(R\) 左 Noether，有限秩自由模的子模有限生成，故核有有限自由满射。重复此论证即可。这里控制的是每个 \(P_n\) 的秩，并没有保证某个关系核 \(K_n\) 变为零，所以每项有限秩并不等于消解只有有限个非零项。\(k[\epsilon]/(\epsilon^2)\) 上的 \(k\) 有逐项秩一的无限消解，下一章还会用它的正次导出函子排除有限长度消解。
\end{solution}

\begin{exercise}\label{b4:ex:07-nonnoetherian-relations}
设 \(k\) 是域，\(R=k[x_1,x_2,\ldots]\)，\(\mathfrak m=(x_1,x_2,\ldots)\)。证明循环模 \(R/\mathfrak m\) 的自然自由满射 \(R\to R/\mathfrak m\) 的核不是有限生成的。
\end{exercise}
\begin{solution}
若 \(\mathfrak m=(f_1,\ldots,f_s)\)，每个 \(f_i\) 只有有限个单项式，每个单项式又只含有限个变量，故有限多个 \(f_i\) 所涉及的变量总并仍有限。取其中未出现的 \(x_j\)，把除 \(x_j\) 外的所有变量置零，得到 \(R\to k[x_j]\) 的同态。各 \(f_i\) 的常数项为零，故其像均为零，而 \(x_j\) 的像非零，与它属于所生成理想矛盾。有限生成对象的关系未必有限生成。
\end{solution}

\begin{exercise}\label{b4:ex:07-cyclic-injective-resolution}
对整数 \(m>1\)，用 \(\mathbb Q/\mathbb Z\) 写出 \(\mathbb Z/m\mathbb Z\) 的两项内射消解，并写清嵌入。
\end{exercise}
\begin{solution}
所需序列是
\[
0\longrightarrow\mathbb Z/m\mathbb Z
\xrightarrow{\bar a\mapsto a/m+\mathbb Z}\mathbb Q/\mathbb Z
\xrightarrow{\times m}\mathbb Q/\mathbb Z\longrightarrow0.
\]
若 \(a'-a=mt\)，则 \(a'/m-a/m=t\in\mathbb Z\)，所以嵌入不依赖剩余类代表。又 \(a/m\in\mathbb Z\) 当且仅当 \(m\mid a\)，故它单射。乘 \(m\) 的核恰为这些分数类；乘 \(m\) 满，因为有理数代表元可以除以 \(m\)。两项可除，因而内射。
\end{solution}

\begin{exercise}\label{b4:ex:07-cyclic-comparison}
设 \(m,n>0\)，\(f:\mathbb Z/m\mathbb Z\to\mathbb Z/n\mathbb Z\) 把 \(\bar1\) 送到 \(\bar c\)。写出它在标准两项自由消解间的提升，并解释 \(c\) 换成 \(c+nt\) 时两个提升的同伦。
\end{exercise}
\begin{solution}
良定义要求 \(n\mid mc\)。零次映射取乘 \(c\)，一次取乘 \(mc/n\)，则 \(n(mc/n)=cm\)，方块交换。记使用 \(c+nt\) 的新提升为 \(u\)，使用 \(c\) 的原提升为 \(v\)。则 \(u-v\) 在零次为乘 \(nt\)，一次为乘 \(mt\)。唯一可能非零的链同伦分量是 \(s_0:P_0\to Q_1\)。取 \(s_0:\mathbb Z\to\mathbb Z\) 为乘 \(t\)，则两个次数的同伦公式分别为
\[
(u-v)_0=\mathrm{d}_1^Qs_0=\times nt,
\qquad
(u-v)_1=s_0\mathrm{d}_1^P=\times mt.
\]
因此它给出所需同伦。
\end{solution}

\begin{exercise}\label{b4:ex:07-projective-source-necessary}
若把比较定理中的“源为投射消解”改为“源为任意正合增广复形”，结论是否仍成立？用 \(\mathbb Z/2\mathbb Z\) 构造反例。
\end{exercise}
\begin{solution}
令源只有零次 \(P_0=\mathbb Z/2\mathbb Z\)，增广为恒等；目标取标准自由消解 \(\mathbb Z\xrightarrow2\mathbb Z\to\mathbb Z/2\mathbb Z\)。若恒等有提升，零次便是 \(\mathbb Z/2\mathbb Z\to\mathbb Z\) 的同态且复合到 \(\mathbb Z/2\mathbb Z\) 为恒等。但任何这样的同态均为零，因为整数无二挠。这排除了提升。失败恰好发生在比较定理的第一步：源的零次项不是投射对象，恒等态射无法沿目标的自由满射提升；即使目标是一份投射消解，也不能弥补源缺少的投射性。
\end{solution}

\begin{exercise}\label{b4:ex:07-add-contractible}
设 \(P_\bullet\to M\) 是投射消解，\(Q\) 是投射对象。在相邻非负次数加入复形 \(Q\xrightarrow1Q\)，其余次数为零。证明所得直和仍是 \(M\) 的投射消解，并写出它到原消解的同伦等价。
\end{exercise}
\begin{solution}
设新增复形 \(T\) 的非零项在次数 \(r+1,r\)，其中 \(r\ge0\)，微分 \(T_{r+1}\to T_r\) 为 \(1_Q\)。在新增零次项上把增广取为零。\(T\) 的同调为零，每次对象仍为投射双积，故增广同调不变。包含 \(i\) 和投影 \(p\) 满足 \(pi=1\)；\(1-ip\) 仅在新增复形上为恒等。取 \(s_r:T_r\to T_{r+1}\) 为 \(1_Q\)，其余同伦分量取零，就有 \(\mathrm{d}s+s\mathrm{d}=1-ip\)。因此 \(i,p\) 是同伦逆。
\end{solution}

\begin{exercise}\label{b4:ex:07-projective-augmented-split}
设 \(M\) 本身投射。证明任何投射消解 \(P_\bullet\to M\) 的增广复形都可缩，并解释这里的收缩为什么一般依赖选择。
\end{exercise}
\begin{solution}
\(\varepsilon:P_0\to M\) 有截面 \(s_{-1}:M\to P_0\)，满足 \(\varepsilon s_{-1}=1_M\)，所以短正合列
\[
0\longrightarrow K_1\longrightarrow P_0\xrightarrow{\varepsilon}M\longrightarrow0
\]
分裂，\(K_1\) 是 \(P_0\) 的投射直和因子。继而
\[
0\longrightarrow K_2\longrightarrow P_1\longrightarrow K_1\longrightarrow0
\]
也分裂，\(K_2\) 因而投射。如此递归，依次得到 \(P_n\cong K_{n+1}\oplus K_n\)，其中 \(K_0=M\)。微分把第二因子恒等送到前一项的第一因子，在另一因子上为零。将各第一因子恒等提升到后一项的第二因子，连同增广复形的次数 \(-1\) 上的 \(s_{-1}:M\to P_0\)，即得 \(\mathrm{d}s+s\mathrm{d}=1\)；在次数 \(-1\) 上，该式就是 \(\varepsilon s_{-1}=1_M\)。截面没有被投射性唯一指定，所以收缩一般依赖选择。
\end{solution}

\begin{exercise}\label{b4:ex:07-horseshoe-p-square}
令 \(p\) 为素数，对 \(0\to\mathbb Z/p\mathbb Z\xrightarrow{\bar a\mapsto p a}\mathbb Z/p^2\mathbb Z\to\mathbb Z/p\mathbb Z\to0\) 构造马蹄消解。验证中间微分可取 \(\begin{pmatrix}p&-1\\0&p\end{pmatrix}\)，增广为 \((a,b)\mapsto pa+b\pmod{p^2}\)。
\end{exercise}
\begin{solution}
增广满射，与微分复合为 \(p^2a\)，故为零；矩阵行列式为 \(p^2\ne0\)，所以它在 \(\mathbb Z^2\) 上单射。若 \(pa+b=p^2t\)，则 \(b=p(pt-a)\)，而
\[
\begin{pmatrix}p&-1\\0&p\end{pmatrix}
\binom{t}{pt-a}=\binom{a}{b}.
\]
因此其像恰为增广核。左端消解以第一坐标包含，右端以第二坐标投影，两个对角块均为乘 \(p\)。非对角项 \(-1\) 使增广复合中的两项 \(-pb\) 与 \(pb\) 抵消，因而把两端消解相容地拼成给定的非分裂扩张，而不是它们的直和。若去掉非对角项 \(-1\)，未增广复形就成为两端消解的直和，零次同调为 \((\mathbb Z/p\mathbb Z)^2\)。而原增广与新微分的复合为 \(p^2a+pb\pmod{p^2}\)，不恒为零，所以它已不再是给定扩张 \(\mathbb Z/p^2\mathbb Z\) 的增广消解。
\end{solution}

\begin{exercise}\label{b4:ex:07-koszul-repeat}
设 \(k\) 是域。在 \(R=k[x]\) 上计算 \(K_\bullet(x,x;R)\) 的全部同调，说明第二个元素的正则性为何失败。
\end{exercise}
\begin{solution}
微分为 \(R\xrightarrow{(-x,x)}R^2\xrightarrow{(x,x)}R\)。左端单射，故 \(H_2=0\)；中间核为 \(R(1,-1)\)，而
\[
\operatorname{im}\mathrm{d}_2=R(-x,x)=Rx(1,-1),
\]
因为乘以单位 \(-1\) 不改变生成的子模。因此
\[
H_1=R(1,-1)/Rx(1,-1)\cong R/(x).
\]\(H_0=R/(x)\)。第二个 \(x\) 在 \(R/(x)\) 上为零，乘法不是单射，正是一次同调不消失的原因。
\end{solution}

\begin{exercise}\label{b4:ex:07-koszul-nonreduced}
设 \(k\) 是域。在 \(R=k[x,y]\) 上写出 \(R/(x,y^2)\) 的两步自由消解。末端商环含有非零幂零元，是否妨碍该消解正合？
\end{exercise}
\begin{solution}
\(x\) 在 \(R\) 上非零因子，\(y^2\) 在 \(R/(x)\cong k[y]\) 上也非零因子。因此消解为
\[
0\to R\xrightarrow{(-y^2,x)}R^2\xrightarrow{(x,y^2)}R\to R/(x,y^2)\to0.
\]
末端 \(\bar y\ne0\) 而 \(\bar y^2=0\)，但正则性是在逐步取商以前检验乘法的单射性，最终商环仍可能非约化。
\end{solution}

\begin{exercise}\label{b4:ex:07-bar-contraction-equivariance}
在\cref{b4:exm:homogeneous-bar}中，假设 \(k\ne0\)，并取非平凡元素 \(g\in G\)。比较 \(s_0(g)\) 与 \(g\cdot s_0(1)\)，并在次数零验证收缩等式。
\end{exercise}
\begin{solution}
作用的两种复合为
\[
s_0\bigl(g\cdot(1)\bigr)=s_0(g)=(1,g),
\qquad
 g\cdot s_0(1)=g\cdot(1,1)=(g,g).
\]
两者是不同基元，故 \(s_0\) 不是 \(G\) 等变态射。又 \(\mathrm{d}_1s_0(g)=(g)-(1)\)，\(s_{-1}\varepsilon(g)=(1)\)，两者之和为 \((g)\)。底层 \(k\) 模上的收缩足以证明正合性，却不能推出平凡 \(kG\) 模投射。
\end{solution}

\begin{exercise}\label{b4:ex:07-finitely-generated-no-injectives}
证明有限生成交换群范畴没有非零内射对象，因而没有足够多内射对象。
\end{exercise}
\begin{solution}
若 \(I\) 在此范畴内射，给定整数 \(n\ge1\) 和 \(a\in I\)，先定义 \(n\mathbb Z\to I\)，使 \(n\mapsto a\)。由单态射 \(n\mathbb Z\hookrightarrow\mathbb Z\) 和内射性，将其延拓为 \(\varphi:\mathbb Z\to I\)。令 \(b=\varphi(1)\)，便有 \(a=\varphi(n)=nb\)，故 \(I\) 可除。有限生成交换群的结构定理给出 \(I\cong\mathbb Z^r\oplus T\)，\(T\) 有限。乘二满射排除 \(r>0\)，乘 \(|T|\) 满射又迫使 \(T=0\)。所以只有零对象内射，而 \(\mathbb Z\) 不能嵌入零对象。
\end{solution}

\begin{exercise}\label{b4:ex:07-cyclic-injective-comparison}
设 \(m,n>1\) 为整数，\(f:\mathbb Z/m\mathbb Z\to\mathbb Z/n\mathbb Z\) 为交换群同态，并设 \(f(\bar1)=\bar c\)，其中 \(c\in\mathbb Z\)。分别用
\[
0\longrightarrow\mathbb Z/r\mathbb Z
\xrightarrow{\bar a\mapsto a/r+\mathbb Z}\mathbb Q/\mathbb Z
\xrightarrow{\times r}\mathbb Q/\mathbb Z\longrightarrow0,
\qquad r=m,n,
\]
作为两者的内射消解。写出 \(f\) 在消解之间的延拓，使两个分量均为乘某个整数的映射。若把 \(c\) 换为 \(c+nt\)（\(t\in\mathbb Z\)），写出两个延拓之间的链同伦。
\end{exercise}
\begin{solution}
同态的良定义要求 \(n\mid mc\)，故 \(\ell=mc/n\) 是整数。记两份内射消解为 \(I_m^\bullet,I_n^\bullet\)。嵌入把 \(\bar1\) 分别送到 \(1/m+\mathbb Z\) 和 \(1/n+\mathbb Z\)。要延拓 \(f\)，零次乘数应使 \(\ell/m+\mathbb Z=c/n+\mathbb Z\)，取 \(\ell=mc/n\) 正好满足这个要求。再由微分相容等式 \(n\ell=mc\) 确定一次分量取乘 \(c\)，所以定义
\[
F^0=\times\ell,\qquad F^1=\times c.
\]
对 \(\bar a\in\mathbb Z/m\mathbb Z\)，有 \(\ell a/m+\mathbb Z=ca/n+\mathbb Z\)，所以 \(F^0\) 延拓 \(f\)。又 \(n\ell=mc\)，故 \((\times n)F^0=F^1(\times m)\)，两个分量组成上链复形映射。

更换代表元后，\(\ell\) 变为 \(\ell+mt\)，两个延拓 \(F',F\) 的差在零次为乘 \(mt\)，在一次为乘 \(nt\)。取唯一可能非零的同伦分量
\[
h^1:I_m^1\longrightarrow I_n^0,\qquad h^1=\times t,
\]
则 \(h^1\mathrm d_{I_m}^0=\times mt\)、\(\mathrm d_{I_n}^0h^1=\times nt\)，正好给出 \(F'-F=\mathrm dh+h\mathrm d\)。
\end{solution}
